\documentclass[]{article}
\usepackage[]{amsmath,amssymb}
\usepackage[]{multirow}
\usepackage{colortbl}
\usepackage[table]{xcolor}
\begin{document}
\title{Assignment 3}
\author{Ketan Sharma}
\date{\today}
\maketitle
\textbf{Que 1}  Write LaTeX code to produce the following expression:
\[
	\overbrace{a+b+c}^{First Group}+\underbrace{x+y+z}_{Second Group}=\overbrace{\underbrace{a+b}_{Sum 1}+c}^{Total}
\]
\textbf{Que 2} Create the following array using the array environment:
	\[
		A=\left[
			\begin{array}{c|ccc}
				 &b&c&d\\ \hline
				a&b&c&d\\
				a&b&c&d\\
				a&b&c&d
			\end{array}
		\right]
	\]
\textbf{Que 3} Write the following piecewise function using the cases environment:	
	\[
		f(x)=\left\{
		\begin{array}{ll}
			x^2+2x+1,&x<0,\\
			\sqrt{x}+1,&0\leq x<4,\\
			\frac{x^2-1}{x-1},&x\geq4	
		\end{array}
		\right.
	\]
\textbf{Que 4}Typeset the following matrix and label its different sections using :
\[
\left[
\begin{array}{ccccc}
	\multicolumn{3}{c}{\overbrace{\phantom{a_{11}\quad a_{12}\quad a_{13}}}^{\text{Column Group 1}}}
	&
	\multicolumn{2}{c}{\overbrace{\phantom{a_{14}\quad a_{15}}}^{\text{Column Group 2}}}
	\\[6pt]
	a_{11} & a_{12} & a_{13} & a_{14} & a_{15}\\
	a_{21} & a_{22} & a_{23} & a_{24} & a_{25}\\
	a_{31} & a_{32} & a_{33} & a_{34} & a_{35}
\end{array}
\right]
\]
\textbf{Que 7} Create the following table.
\begin{center}
	\begin{tabular}{|c|c|c|}
		\hline
		\multicolumn{3}{|c|}{Academic Performance} \\ \hline
			a&b&c\\
			1&2&3\\
			1&2&3\\ \hline
	\end{tabular}
\end{center}
\textbf{Que 8}Create the following table.
\begin{center}
	\begin{tabular}{|c|l|c|c|}
		\hline
		Department&Course&Credits&Grade\\
		\hline
		\multirow{3}{*}{CSE}
		&Data Structures & 4&A+\\
		&Machine Learning&4&A+\\
		&Operating Systems&3&B+\\ \hline
		\multirow{2}{*}{ECE}
		&Digital Electronics &4&A\\
		&Signals&3&B+\\ \hline
	\end{tabular}
\end{center}
\textbf{Que 9}Create the following scientific table.
\textbf{Que 9} Create the following scientific table.
\begin{center}
	\begin{tabular}{|c|l|c|c|}
		\hline
		\rowcolor{blue}
		Department&Course&Credits&Grade\\
		\hline
		
		\cellcolor{white}
		\multirow{3}{*}{CSE}
		&Data Structures & 4&A+\\
		
		\rowcolor{cyan}
		&Machine Learning&4&A+\\
		
		&Operating Systems&3&B+\\
		\hline
		
		\cellcolor{white}
		\multirow{2}{*}{ECE}
		&Digital Electronics &4&A\\
		
		\rowcolor{cyan}
		&Signals&3&B+\\
		\hline
	\end{tabular}
\end{center}
\noindent
\textbf{Que 10} Create a student-performance table.

\begin{center}
	\begin{tabular}{|c|c|c|c|}
		\hline
		\rowcolor{blue}
		Student & Mathematics & Physics & Computer Science\\
		\hline
		
		Amit & \cellcolor{green!20}92 & \cellcolor{green!20}81 & \cellcolor{green!20}95\\
		\hline
		Neha & \cellcolor{red!20}72 & \cellcolor{yellow!30}88 & \cellcolor{green!20}91\\
		\hline
		Rahul & \cellcolor{red!20}68 & \cellcolor{red!20}74 & \cellcolor{yellow!30}79\\
		\hline
		Priya & \cellcolor{green!20}96 & \cellcolor{green!20}94 & \cellcolor{green!20}98\\
		\hline
	\end{tabular}
\end{center}


\end{document}